\section{账本事件的区域史深描：northern-song}

\label{sec:ledger-depth-northern-song}

本组把账本事件置于战争、制度、区域社会与证据类型的交叉处。年表提供政治节点，地方材料与物质遗存则限制我们对征发、身份、信仰和迁徙所能作出的断言。

\subsection{战争、行政与区域资源}

账本条目\texttt{NS-0960-CHENQIAO}所记960年正月的“陈桥兵变后赵匡胤即帝位，改国号宋”，属于后周禁军与赵匡胤集团的历史语境。条目所示前因是后周幼主继位、禁军将领掌兵与北伐动员共同造成权力转换条件，其可见后果为宋廷取得汴京中枢和后周官僚资源，但须继续处理藩镇、割据政权与契丹边患。这一变化若进入区域生活，必须通过道路、文书、军队、市场、家庭和地方中介才会发生；政权易手或法令发布不等于所有城镇、乡村与边地在同一时间获得同样经验。

本条所列来源为《宋史》卷1〈太祖本纪一〉；《续资治通鉴长编》建隆元年正月条；《宋会要辑稿·礼》建国礼仪条。它们能够支持特定的年序、制度语言、政治行动或局部遗存，却不能无条件化为社会总量。账本特别提醒“核心事实可靠；黄袍加身仪式细节多出自后起叙事”。因此，军报的兵数、条约的名分、官府的族群分类、判牍中的道德判断以及后出的叙述都须回到产生它们的场景，而不可直接代表全域动机或人口。

将事件与地方层次连接时，应同时考察资源怎样被征集、信息怎样传递、迁徙怎样改变户籍和劳动、宗教与亲属网络怎样提供或限制协作。现有证据往往只照亮少数地点；保留这种不均衡，才能避免把宋辽夏金元的多政权历史改写成单一中心的必然过程。

账本条目\texttt{NS-0960-KAIFENG-CAPITAL}所记960年的“宋廷承用开封为东京和中央行政中心”，属于宋的历史语境。条目所示前因是新朝须占有后周既有都城、仓储和运河节点以维持禁军与百官，其可见后果为以汴河—开封为中心的财政与交通结构成为统一战争和城市供给的长期条件。这一变化若进入区域生活，必须通过道路、文书、军队、市场、家庭和地方中介才会发生；政权易手或法令发布不等于所有城镇、乡村与边地在同一时间获得同样经验。

本条所列来源为《宋史》卷1〈太祖本纪一〉；《宋会要辑稿·方域》京师条；开封北宋东京城遗址城垣、街道与水系考古资料。它们能够支持特定的年序、制度语言、政治行动或局部遗存，却不能无条件化为社会总量。账本特别提醒“建都事实可靠；城坊、水系的具体承袭层次须由考古分期核对”。因此，军报的兵数、条约的名分、官府的族群分类、判牍中的道德判断以及后出的叙述都须回到产生它们的场景，而不可直接代表全域动机或人口。

将事件与地方层次连接时，应同时考察资源怎样被征集、信息怎样传递、迁徙怎样改变户籍和劳动、宗教与亲属网络怎样提供或限制协作。现有证据往往只照亮少数地点；保留这种不均衡，才能避免把宋辽夏金元的多政权历史改写成单一中心的必然过程。

账本条目\texttt{NS-0960-REGNAL-YEAR}所记960年的“宋以宋太祖在位第1年作为诏令、赋役与外交文书的纪年框架”，属于宋的历史语境。条目所示前因是新岁颁历、年号和纪年是中枢与地方行政沟通的共同前提，其可见后果为为同年政令、战事和地方材料提供日期锚点，不能单独推知具体政策执行。这一变化若进入区域生活，必须通过道路、文书、军队、市场、家庭和地方中介才会发生；政权易手或法令发布不等于所有城镇、乡村与边地在同一时间获得同样经验。

本条所列来源为《宋史》卷1—23〈本纪〉；《续资治通鉴长编》本年条；《宋会要辑稿》相关门类。它们能够支持特定的年序、制度语言、政治行动或局部遗存，却不能无条件化为社会总量。账本特别提醒“纪年框架可靠；该记录仅确证政权纪年连续，年内具体事项须由本年条另证”。因此，军报的兵数、条约的名分、官府的族群分类、判牍中的道德判断以及后出的叙述都须回到产生它们的场景，而不可直接代表全域动机或人口。

将事件与地方层次连接时，应同时考察资源怎样被征集、信息怎样传递、迁徙怎样改变户籍和劳动、宗教与亲属网络怎样提供或限制协作。现有证据往往只照亮少数地点；保留这种不均衡，才能避免把宋辽夏金元的多政权历史改写成单一中心的必然过程。

账本条目\texttt{NS-0961-MILITARY-POWER}所记961年前后的“太祖以宴饮、任官和调动方式削弱禁军宿将的独立兵权”，属于宋的历史语境。条目所示前因是五代以来统兵将领可凭禁军发动政变，新朝需降低同类风险，其可见后果为中央以轮调、文官监军和财赋供养重塑军事指挥，亦增加对东京财政的依赖。这一变化若进入区域生活，必须通过道路、文书、军队、市场、家庭和地方中介才会发生；政权易手或法令发布不等于所有城镇、乡村与边地在同一时间获得同样经验。

本条所列来源为《宋史》卷1〈太祖本纪一〉；《宋史》卷250〈石守信传〉；《续资治通鉴长编》建隆二年条。它们能够支持特定的年序、制度语言、政治行动或局部遗存，却不能无条件化为社会总量。账本特别提醒“权力重组可靠；后世所谓杯酒释兵权的对话与一次性发生说有文学化可能”。因此，军报的兵数、条约的名分、官府的族群分类、判牍中的道德判断以及后出的叙述都须回到产生它们的场景，而不可直接代表全域动机或人口。

将事件与地方层次连接时，应同时考察资源怎样被征集、信息怎样传递、迁徙怎样改变户籍和劳动、宗教与亲属网络怎样提供或限制协作。现有证据往往只照亮少数地点；保留这种不均衡，才能避免把宋辽夏金元的多政权历史改写成单一中心的必然过程。

账本条目\texttt{NS-0961-REGNAL-YEAR}所记961年的“宋以宋太祖在位第2年作为诏令、赋役与外交文书的纪年框架”，属于宋的历史语境。条目所示前因是新岁颁历、年号和纪年是中枢与地方行政沟通的共同前提，其可见后果为为同年政令、战事和地方材料提供日期锚点，不能单独推知具体政策执行。这一变化若进入区域生活，必须通过道路、文书、军队、市场、家庭和地方中介才会发生；政权易手或法令发布不等于所有城镇、乡村与边地在同一时间获得同样经验。

本条所列来源为《宋史》卷1—23〈本纪〉；《续资治通鉴长编》本年条；《宋会要辑稿》相关门类。它们能够支持特定的年序、制度语言、政治行动或局部遗存，却不能无条件化为社会总量。账本特别提醒“纪年框架可靠；该记录仅确证政权纪年连续，年内具体事项须由本年条另证”。因此，军报的兵数、条约的名分、官府的族群分类、判牍中的道德判断以及后出的叙述都须回到产生它们的场景，而不可直接代表全域动机或人口。

将事件与地方层次连接时，应同时考察资源怎样被征集、信息怎样传递、迁徙怎样改变户籍和劳动、宗教与亲属网络怎样提供或限制协作。现有证据往往只照亮少数地点；保留这种不均衡，才能避免把宋辽夏金元的多政权历史改写成单一中心的必然过程。

账本条目\texttt{NS-0962-REGNAL-YEAR}所记962年的“宋以宋太祖在位第3年作为诏令、赋役与外交文书的纪年框架”，属于宋的历史语境。条目所示前因是新岁颁历、年号和纪年是中枢与地方行政沟通的共同前提，其可见后果为为同年政令、战事和地方材料提供日期锚点，不能单独推知具体政策执行。这一变化若进入区域生活，必须通过道路、文书、军队、市场、家庭和地方中介才会发生；政权易手或法令发布不等于所有城镇、乡村与边地在同一时间获得同样经验。

本条所列来源为《宋史》卷1—23〈本纪〉；《续资治通鉴长编》本年条；《宋会要辑稿》相关门类。它们能够支持特定的年序、制度语言、政治行动或局部遗存，却不能无条件化为社会总量。账本特别提醒“纪年框架可靠；该记录仅确证政权纪年连续，年内具体事项须由本年条另证”。因此，军报的兵数、条约的名分、官府的族群分类、判牍中的道德判断以及后出的叙述都须回到产生它们的场景，而不可直接代表全域动机或人口。

将事件与地方层次连接时，应同时考察资源怎样被征集、信息怎样传递、迁徙怎样改变户籍和劳动、宗教与亲属网络怎样提供或限制协作。现有证据往往只照亮少数地点；保留这种不均衡，才能避免把宋辽夏金元的多政权历史改写成单一中心的必然过程。

账本条目\texttt{NS-0963-JINGNAN}所记963年的“宋军入荆南，结束高氏政权”，属于宋与荆南的历史语境。条目所示前因是荆南体量小且夹处宋、后蜀、南唐之间，难以保持独立军政，其可见后果为宋取得长江中游战略节点，并以改置州县试验对旧国行政资源的接收。这一变化若进入区域生活，必须通过道路、文书、军队、市场、家庭和地方中介才会发生；政权易手或法令发布不等于所有城镇、乡村与边地在同一时间获得同样经验。

本条所列来源为《宋史》卷2〈太祖本纪二〉；《续资治通鉴长编》乾德元年条；《宋会要辑稿·方域》荆南条。它们能够支持特定的年序、制度语言、政治行动或局部遗存，却不能无条件化为社会总量。账本特别提醒“年代可靠；投降过程与军队规模须谨慎采用宋方叙事”。因此，军报的兵数、条约的名分、官府的族群分类、判牍中的道德判断以及后出的叙述都须回到产生它们的场景，而不可直接代表全域动机或人口。

将事件与地方层次连接时，应同时考察资源怎样被征集、信息怎样传递、迁徙怎样改变户籍和劳动、宗教与亲属网络怎样提供或限制协作。现有证据往往只照亮少数地点；保留这种不均衡，才能避免把宋辽夏金元的多政权历史改写成单一中心的必然过程。

账本条目\texttt{NS-0963-REGNAL-YEAR}所记963年的“宋以宋太祖在位第4年作为诏令、赋役与外交文书的纪年框架”，属于宋的历史语境。条目所示前因是新岁颁历、年号和纪年是中枢与地方行政沟通的共同前提，其可见后果为为同年政令、战事和地方材料提供日期锚点，不能单独推知具体政策执行。这一变化若进入区域生活，必须通过道路、文书、军队、市场、家庭和地方中介才会发生；政权易手或法令发布不等于所有城镇、乡村与边地在同一时间获得同样经验。

本条所列来源为《宋史》卷1—23〈本纪〉；《续资治通鉴长编》本年条；《宋会要辑稿》相关门类。它们能够支持特定的年序、制度语言、政治行动或局部遗存，却不能无条件化为社会总量。账本特别提醒“纪年框架可靠；该记录仅确证政权纪年连续，年内具体事项须由本年条另证”。因此，军报的兵数、条约的名分、官府的族群分类、判牍中的道德判断以及后出的叙述都须回到产生它们的场景，而不可直接代表全域动机或人口。

将事件与地方层次连接时，应同时考察资源怎样被征集、信息怎样传递、迁徙怎样改变户籍和劳动、宗教与亲属网络怎样提供或限制协作。现有证据往往只照亮少数地点；保留这种不均衡，才能避免把宋辽夏金元的多政权历史改写成单一中心的必然过程。

账本条目\texttt{NS-0964-LATER-SHU-CAMPAIGN}所记964年冬的“宋军发动伐后蜀战役”，属于宋与后蜀的历史语境。条目所示前因是宋先稳住荆南，意在消除西南割据并取得成都平原资源，其可见后果为战争打破后蜀防御体系，宋军开始面对山地运输、降附军民安置和四川财政整合。这一变化若进入区域生活，必须通过道路、文书、军队、市场、家庭和地方中介才会发生；政权易手或法令发布不等于所有城镇、乡村与边地在同一时间获得同样经验。

本条所列来源为《宋史》卷2〈太祖本纪二〉；《宋史》卷479〈后蜀世家〉；《续资治通鉴长编》乾德二年冬条。它们能够支持特定的年序、制度语言、政治行动或局部遗存，却不能无条件化为社会总量。账本特别提醒“出兵日期可靠；行军路线和战功归属有异说”。因此，军报的兵数、条约的名分、官府的族群分类、判牍中的道德判断以及后出的叙述都须回到产生它们的场景，而不可直接代表全域动机或人口。

将事件与地方层次连接时，应同时考察资源怎样被征集、信息怎样传递、迁徙怎样改变户籍和劳动、宗教与亲属网络怎样提供或限制协作。现有证据往往只照亮少数地点；保留这种不均衡，才能避免把宋辽夏金元的多政权历史改写成单一中心的必然过程。

账本条目\texttt{NS-0964-REGNAL-YEAR}所记964年的“宋以宋太祖在位第5年作为诏令、赋役与外交文书的纪年框架”，属于宋的历史语境。条目所示前因是新岁颁历、年号和纪年是中枢与地方行政沟通的共同前提，其可见后果为为同年政令、战事和地方材料提供日期锚点，不能单独推知具体政策执行。这一变化若进入区域生活，必须通过道路、文书、军队、市场、家庭和地方中介才会发生；政权易手或法令发布不等于所有城镇、乡村与边地在同一时间获得同样经验。

本条所列来源为《宋史》卷1—23〈本纪〉；《续资治通鉴长编》本年条；《宋会要辑稿》相关门类。它们能够支持特定的年序、制度语言、政治行动或局部遗存，却不能无条件化为社会总量。账本特别提醒“纪年框架可靠；该记录仅确证政权纪年连续，年内具体事项须由本年条另证”。因此，军报的兵数、条约的名分、官府的族群分类、判牍中的道德判断以及后出的叙述都须回到产生它们的场景，而不可直接代表全域动机或人口。

将事件与地方层次连接时，应同时考察资源怎样被征集、信息怎样传递、迁徙怎样改变户籍和劳动、宗教与亲属网络怎样提供或限制协作。现有证据往往只照亮少数地点；保留这种不均衡，才能避免把宋辽夏金元的多政权历史改写成单一中心的必然过程。

账本条目\texttt{NS-0965-LATER-SHU-ANNEXATION}所记965年的“孟昶降宋，后蜀政权终结”，属于宋与后蜀的历史语境。条目所示前因是成都失守后后蜀朝廷无力继续组织防御，其可见后果为宋获得西南富庶区，却因军纪、赋役与旧官安置问题埋下蜀地反抗条件。这一变化若进入区域生活，必须通过道路、文书、军队、市场、家庭和地方中介才会发生；政权易手或法令发布不等于所有城镇、乡村与边地在同一时间获得同样经验。

本条所列来源为《宋史》卷2〈太祖本纪二〉；《宋史》卷479〈后蜀世家〉；《续资治通鉴长编》乾德三年条。它们能够支持特定的年序、制度语言、政治行动或局部遗存，却不能无条件化为社会总量。账本特别提醒“降服事实可靠；宋军入蜀后的掠夺与赏赐范围须分证”。因此，军报的兵数、条约的名分、官府的族群分类、判牍中的道德判断以及后出的叙述都须回到产生它们的场景，而不可直接代表全域动机或人口。

将事件与地方层次连接时，应同时考察资源怎样被征集、信息怎样传递、迁徙怎样改变户籍和劳动、宗教与亲属网络怎样提供或限制协作。现有证据往往只照亮少数地点；保留这种不均衡，才能避免把宋辽夏金元的多政权历史改写成单一中心的必然过程。

账本条目\texttt{NS-0965-REGNAL-YEAR}所记965年的“宋以宋太祖在位第6年作为诏令、赋役与外交文书的纪年框架”，属于宋的历史语境。条目所示前因是新岁颁历、年号和纪年是中枢与地方行政沟通的共同前提，其可见后果为为同年政令、战事和地方材料提供日期锚点，不能单独推知具体政策执行。这一变化若进入区域生活，必须通过道路、文书、军队、市场、家庭和地方中介才会发生；政权易手或法令发布不等于所有城镇、乡村与边地在同一时间获得同样经验。

本条所列来源为《宋史》卷1—23〈本纪〉；《续资治通鉴长编》本年条；《宋会要辑稿》相关门类。它们能够支持特定的年序、制度语言、政治行动或局部遗存，却不能无条件化为社会总量。账本特别提醒“纪年框架可靠；该记录仅确证政权纪年连续，年内具体事项须由本年条另证”。因此，军报的兵数、条约的名分、官府的族群分类、判牍中的道德判断以及后出的叙述都须回到产生它们的场景，而不可直接代表全域动机或人口。

将事件与地方层次连接时，应同时考察资源怎样被征集、信息怎样传递、迁徙怎样改变户籍和劳动、宗教与亲属网络怎样提供或限制协作。现有证据往往只照亮少数地点；保留这种不均衡，才能避免把宋辽夏金元的多政权历史改写成单一中心的必然过程。

\subsection{外交、道路与人口移动}

账本条目\texttt{NS-0966-REGNAL-YEAR}所记966年的“宋以宋太祖在位第7年作为诏令、赋役与外交文书的纪年框架”，属于宋的历史语境。条目所示前因是新岁颁历、年号和纪年是中枢与地方行政沟通的共同前提，其可见后果为为同年政令、战事和地方材料提供日期锚点，不能单独推知具体政策执行。这一变化若进入区域生活，必须通过道路、文书、军队、市场、家庭和地方中介才会发生；政权易手或法令发布不等于所有城镇、乡村与边地在同一时间获得同样经验。

本条所列来源为《宋史》卷1—23〈本纪〉；《续资治通鉴长编》本年条；《宋会要辑稿》相关门类。它们能够支持特定的年序、制度语言、政治行动或局部遗存，却不能无条件化为社会总量。账本特别提醒“纪年框架可靠；该记录仅确证政权纪年连续，年内具体事项须由本年条另证”。因此，军报的兵数、条约的名分、官府的族群分类、判牍中的道德判断以及后出的叙述都须回到产生它们的场景，而不可直接代表全域动机或人口。

将事件与地方层次连接时，应同时考察资源怎样被征集、信息怎样传递、迁徙怎样改变户籍和劳动、宗教与亲属网络怎样提供或限制协作。现有证据往往只照亮少数地点；保留这种不均衡，才能避免把宋辽夏金元的多政权历史改写成单一中心的必然过程。

账本条目\texttt{NS-0967-REGNAL-YEAR}所记967年的“宋以宋太祖在位第8年作为诏令、赋役与外交文书的纪年框架”，属于宋的历史语境。条目所示前因是新岁颁历、年号和纪年是中枢与地方行政沟通的共同前提，其可见后果为为同年政令、战事和地方材料提供日期锚点，不能单独推知具体政策执行。这一变化若进入区域生活，必须通过道路、文书、军队、市场、家庭和地方中介才会发生；政权易手或法令发布不等于所有城镇、乡村与边地在同一时间获得同样经验。

本条所列来源为《宋史》卷1—23〈本纪〉；《续资治通鉴长编》本年条；《宋会要辑稿》相关门类。它们能够支持特定的年序、制度语言、政治行动或局部遗存，却不能无条件化为社会总量。账本特别提醒“纪年框架可靠；该记录仅确证政权纪年连续，年内具体事项须由本年条另证”。因此，军报的兵数、条约的名分、官府的族群分类、判牍中的道德判断以及后出的叙述都须回到产生它们的场景，而不可直接代表全域动机或人口。

将事件与地方层次连接时，应同时考察资源怎样被征集、信息怎样传递、迁徙怎样改变户籍和劳动、宗教与亲属网络怎样提供或限制协作。现有证据往往只照亮少数地点；保留这种不均衡，才能避免把宋辽夏金元的多政权历史改写成单一中心的必然过程。

账本条目\texttt{NS-0968-REGNAL-YEAR}所记968年的“宋以宋太祖在位第9年作为诏令、赋役与外交文书的纪年框架”，属于宋的历史语境。条目所示前因是新岁颁历、年号和纪年是中枢与地方行政沟通的共同前提，其可见后果为为同年政令、战事和地方材料提供日期锚点，不能单独推知具体政策执行。这一变化若进入区域生活，必须通过道路、文书、军队、市场、家庭和地方中介才会发生；政权易手或法令发布不等于所有城镇、乡村与边地在同一时间获得同样经验。

本条所列来源为《宋史》卷1—23〈本纪〉；《续资治通鉴长编》本年条；《宋会要辑稿》相关门类。它们能够支持特定的年序、制度语言、政治行动或局部遗存，却不能无条件化为社会总量。账本特别提醒“纪年框架可靠；该记录仅确证政权纪年连续，年内具体事项须由本年条另证”。因此，军报的兵数、条约的名分、官府的族群分类、判牍中的道德判断以及后出的叙述都须回到产生它们的场景，而不可直接代表全域动机或人口。

将事件与地方层次连接时，应同时考察资源怎样被征集、信息怎样传递、迁徙怎样改变户籍和劳动、宗教与亲属网络怎样提供或限制协作。现有证据往往只照亮少数地点；保留这种不均衡，才能避免把宋辽夏金元的多政权历史改写成单一中心的必然过程。

账本条目\texttt{NS-0969-REGNAL-YEAR}所记969年的“宋以宋太祖在位第10年作为诏令、赋役与外交文书的纪年框架”，属于宋的历史语境。条目所示前因是新岁颁历、年号和纪年是中枢与地方行政沟通的共同前提，其可见后果为为同年政令、战事和地方材料提供日期锚点，不能单独推知具体政策执行。这一变化若进入区域生活，必须通过道路、文书、军队、市场、家庭和地方中介才会发生；政权易手或法令发布不等于所有城镇、乡村与边地在同一时间获得同样经验。

本条所列来源为《宋史》卷1—23〈本纪〉；《续资治通鉴长编》本年条；《宋会要辑稿》相关门类。它们能够支持特定的年序、制度语言、政治行动或局部遗存，却不能无条件化为社会总量。账本特别提醒“纪年框架可靠；该记录仅确证政权纪年连续，年内具体事项须由本年条另证”。因此，军报的兵数、条约的名分、官府的族群分类、判牍中的道德判断以及后出的叙述都须回到产生它们的场景，而不可直接代表全域动机或人口。

将事件与地方层次连接时，应同时考察资源怎样被征集、信息怎样传递、迁徙怎样改变户籍和劳动、宗教与亲属网络怎样提供或限制协作。现有证据往往只照亮少数地点；保留这种不均衡，才能避免把宋辽夏金元的多政权历史改写成单一中心的必然过程。

账本条目\texttt{NS-0970-REGNAL-YEAR}所记970年的“宋以宋太祖在位第11年作为诏令、赋役与外交文书的纪年框架”，属于宋的历史语境。条目所示前因是新岁颁历、年号和纪年是中枢与地方行政沟通的共同前提，其可见后果为为同年政令、战事和地方材料提供日期锚点，不能单独推知具体政策执行。这一变化若进入区域生活，必须通过道路、文书、军队、市场、家庭和地方中介才会发生；政权易手或法令发布不等于所有城镇、乡村与边地在同一时间获得同样经验。

本条所列来源为《宋史》卷1—23〈本纪〉；《续资治通鉴长编》本年条；《宋会要辑稿》相关门类。它们能够支持特定的年序、制度语言、政治行动或局部遗存，却不能无条件化为社会总量。账本特别提醒“纪年框架可靠；该记录仅确证政权纪年连续，年内具体事项须由本年条另证”。因此，军报的兵数、条约的名分、官府的族群分类、判牍中的道德判断以及后出的叙述都须回到产生它们的场景，而不可直接代表全域动机或人口。

将事件与地方层次连接时，应同时考察资源怎样被征集、信息怎样传递、迁徙怎样改变户籍和劳动、宗教与亲属网络怎样提供或限制协作。现有证据往往只照亮少数地点；保留这种不均衡，才能避免把宋辽夏金元的多政权历史改写成单一中心的必然过程。

账本条目\texttt{NS-0971-REGNAL-YEAR}所记971年的“宋以宋太祖在位第12年作为诏令、赋役与外交文书的纪年框架”，属于宋的历史语境。条目所示前因是新岁颁历、年号和纪年是中枢与地方行政沟通的共同前提，其可见后果为为同年政令、战事和地方材料提供日期锚点，不能单独推知具体政策执行。这一变化若进入区域生活，必须通过道路、文书、军队、市场、家庭和地方中介才会发生；政权易手或法令发布不等于所有城镇、乡村与边地在同一时间获得同样经验。

本条所列来源为《宋史》卷1—23〈本纪〉；《续资治通鉴长编》本年条；《宋会要辑稿》相关门类。它们能够支持特定的年序、制度语言、政治行动或局部遗存，却不能无条件化为社会总量。账本特别提醒“纪年框架可靠；该记录仅确证政权纪年连续，年内具体事项须由本年条另证”。因此，军报的兵数、条约的名分、官府的族群分类、判牍中的道德判断以及后出的叙述都须回到产生它们的场景，而不可直接代表全域动机或人口。

将事件与地方层次连接时，应同时考察资源怎样被征集、信息怎样传递、迁徙怎样改变户籍和劳动、宗教与亲属网络怎样提供或限制协作。现有证据往往只照亮少数地点；保留这种不均衡，才能避免把宋辽夏金元的多政权历史改写成单一中心的必然过程。

账本条目\texttt{NS-0971-SOUTHERN-HAN}所记971年的“宋军攻取广州，南汉刘鋹降服”，属于宋与南汉的历史语境。条目所示前因是宋为完成南方统一并控制岭南财赋、港口和盐业而南进，其可见后果为岭南纳入宋朝路州体系，朝廷需调适边地军政和既有海上贸易网络。这一变化若进入区域生活，必须通过道路、文书、军队、市场、家庭和地方中介才会发生；政权易手或法令发布不等于所有城镇、乡村与边地在同一时间获得同样经验。

本条所列来源为《宋史》卷2〈太祖本纪二〉；《宋史》卷480〈南汉世家〉；《续资治通鉴长编》开宝四年条。它们能够支持特定的年序、制度语言、政治行动或局部遗存，却不能无条件化为社会总量。账本特别提醒“核心过程可靠；伤亡、俘获与海上封锁数字多为战报”。因此，军报的兵数、条约的名分、官府的族群分类、判牍中的道德判断以及后出的叙述都须回到产生它们的场景，而不可直接代表全域动机或人口。

将事件与地方层次连接时，应同时考察资源怎样被征集、信息怎样传递、迁徙怎样改变户籍和劳动、宗教与亲属网络怎样提供或限制协作。现有证据往往只照亮少数地点；保留这种不均衡，才能避免把宋辽夏金元的多政权历史改写成单一中心的必然过程。

账本条目\texttt{NS-0972-REGNAL-YEAR}所记972年的“宋以宋太祖在位第13年作为诏令、赋役与外交文书的纪年框架”，属于宋的历史语境。条目所示前因是新岁颁历、年号和纪年是中枢与地方行政沟通的共同前提，其可见后果为为同年政令、战事和地方材料提供日期锚点，不能单独推知具体政策执行。这一变化若进入区域生活，必须通过道路、文书、军队、市场、家庭和地方中介才会发生；政权易手或法令发布不等于所有城镇、乡村与边地在同一时间获得同样经验。

本条所列来源为《宋史》卷1—23〈本纪〉；《续资治通鉴长编》本年条；《宋会要辑稿》相关门类。它们能够支持特定的年序、制度语言、政治行动或局部遗存，却不能无条件化为社会总量。账本特别提醒“纪年框架可靠；该记录仅确证政权纪年连续，年内具体事项须由本年条另证”。因此，军报的兵数、条约的名分、官府的族群分类、判牍中的道德判断以及后出的叙述都须回到产生它们的场景，而不可直接代表全域动机或人口。

将事件与地方层次连接时，应同时考察资源怎样被征集、信息怎样传递、迁徙怎样改变户籍和劳动、宗教与亲属网络怎样提供或限制协作。现有证据往往只照亮少数地点；保留这种不均衡，才能避免把宋辽夏金元的多政权历史改写成单一中心的必然过程。

账本条目\texttt{NS-0973-REGNAL-YEAR}所记973年的“宋以宋太祖在位第14年作为诏令、赋役与外交文书的纪年框架”，属于宋的历史语境。条目所示前因是新岁颁历、年号和纪年是中枢与地方行政沟通的共同前提，其可见后果为为同年政令、战事和地方材料提供日期锚点，不能单独推知具体政策执行。这一变化若进入区域生活，必须通过道路、文书、军队、市场、家庭和地方中介才会发生；政权易手或法令发布不等于所有城镇、乡村与边地在同一时间获得同样经验。

本条所列来源为《宋史》卷1—23〈本纪〉；《续资治通鉴长编》本年条；《宋会要辑稿》相关门类。它们能够支持特定的年序、制度语言、政治行动或局部遗存，却不能无条件化为社会总量。账本特别提醒“纪年框架可靠；该记录仅确证政权纪年连续，年内具体事项须由本年条另证”。因此，军报的兵数、条约的名分、官府的族群分类、判牍中的道德判断以及后出的叙述都须回到产生它们的场景，而不可直接代表全域动机或人口。

将事件与地方层次连接时，应同时考察资源怎样被征集、信息怎样传递、迁徙怎样改变户籍和劳动、宗教与亲属网络怎样提供或限制协作。现有证据往往只照亮少数地点；保留这种不均衡，才能避免把宋辽夏金元的多政权历史改写成单一中心的必然过程。

账本条目\texttt{NS-0974-REGNAL-YEAR}所记974年的“宋以宋太祖在位第15年作为诏令、赋役与外交文书的纪年框架”，属于宋的历史语境。条目所示前因是新岁颁历、年号和纪年是中枢与地方行政沟通的共同前提，其可见后果为为同年政令、战事和地方材料提供日期锚点，不能单独推知具体政策执行。这一变化若进入区域生活，必须通过道路、文书、军队、市场、家庭和地方中介才会发生；政权易手或法令发布不等于所有城镇、乡村与边地在同一时间获得同样经验。

本条所列来源为《宋史》卷1—23〈本纪〉；《续资治通鉴长编》本年条；《宋会要辑稿》相关门类。它们能够支持特定的年序、制度语言、政治行动或局部遗存，却不能无条件化为社会总量。账本特别提醒“纪年框架可靠；该记录仅确证政权纪年连续，年内具体事项须由本年条另证”。因此，军报的兵数、条约的名分、官府的族群分类、判牍中的道德判断以及后出的叙述都须回到产生它们的场景，而不可直接代表全域动机或人口。

将事件与地方层次连接时，应同时考察资源怎样被征集、信息怎样传递、迁徙怎样改变户籍和劳动、宗教与亲属网络怎样提供或限制协作。现有证据往往只照亮少数地点；保留这种不均衡，才能避免把宋辽夏金元的多政权历史改写成单一中心的必然过程。

账本条目\texttt{NS-0974-SOUTHERN-TANG-CAMPAIGN}所记974年秋的“宋军渡江围攻金陵，展开灭南唐战争”，属于宋与南唐的历史语境。条目所示前因是南唐虽称臣仍保有江南政权和长江防线，宋欲完成东南统一，其可见后果为江南战争将宋的财政、漕运和攻城能力集中到长江下游。这一变化若进入区域生活，必须通过道路、文书、军队、市场、家庭和地方中介才会发生；政权易手或法令发布不等于所有城镇、乡村与边地在同一时间获得同样经验。

本条所列来源为《宋史》卷2〈太祖本纪二〉；《宋史》卷478〈南唐世家〉；《续资治通鉴长编》开宝七年条。它们能够支持特定的年序、制度语言、政治行动或局部遗存，却不能无条件化为社会总量。账本特别提醒“战争爆发可靠；兵力与战场细节须与不同传记互校”。因此，军报的兵数、条约的名分、官府的族群分类、判牍中的道德判断以及后出的叙述都须回到产生它们的场景，而不可直接代表全域动机或人口。

将事件与地方层次连接时，应同时考察资源怎样被征集、信息怎样传递、迁徙怎样改变户籍和劳动、宗教与亲属网络怎样提供或限制协作。现有证据往往只照亮少数地点；保留这种不均衡，才能避免把宋辽夏金元的多政权历史改写成单一中心的必然过程。

账本条目\texttt{NS-0975-REGNAL-YEAR}所记975年的“宋以宋太祖在位第16年作为诏令、赋役与外交文书的纪年框架”，属于宋的历史语境。条目所示前因是新岁颁历、年号和纪年是中枢与地方行政沟通的共同前提，其可见后果为为同年政令、战事和地方材料提供日期锚点，不能单独推知具体政策执行。这一变化若进入区域生活，必须通过道路、文书、军队、市场、家庭和地方中介才会发生；政权易手或法令发布不等于所有城镇、乡村与边地在同一时间获得同样经验。

本条所列来源为《宋史》卷1—23〈本纪〉；《续资治通鉴长编》本年条；《宋会要辑稿》相关门类。它们能够支持特定的年序、制度语言、政治行动或局部遗存，却不能无条件化为社会总量。账本特别提醒“纪年框架可靠；该记录仅确证政权纪年连续，年内具体事项须由本年条另证”。因此，军报的兵数、条约的名分、官府的族群分类、判牍中的道德判断以及后出的叙述都须回到产生它们的场景，而不可直接代表全域动机或人口。

将事件与地方层次连接时，应同时考察资源怎样被征集、信息怎样传递、迁徙怎样改变户籍和劳动、宗教与亲属网络怎样提供或限制协作。现有证据往往只照亮少数地点；保留这种不均衡，才能避免把宋辽夏金元的多政权历史改写成单一中心的必然过程。

\subsection{环境风险与地方经济}

账本条目\texttt{NS-0975-SOUTHERN-TANG-ANNEXATION}所记975年十一月的“金陵陷落，李煜降宋，南唐灭亡”，属于宋与南唐的历史语境。条目所示前因是金陵久围而援军不至，南唐军事和财政资源耗尽，其可见后果为宋控制江南税粮、手工业和水网地区，统一重心转向北汉与辽边。这一变化若进入区域生活，必须通过道路、文书、军队、市场、家庭和地方中介才会发生；政权易手或法令发布不等于所有城镇、乡村与边地在同一时间获得同样经验。

本条所列来源为《宋史》卷2〈太祖本纪二〉；《宋史》卷478〈南唐世家〉；《续资治通鉴长编》开宝八年十一月条。它们能够支持特定的年序、制度语言、政治行动或局部遗存，却不能无条件化为社会总量。账本特别提醒“可靠纪年；李煜待遇和遗民反应不应以文学作品反推全体社会”。因此，军报的兵数、条约的名分、官府的族群分类、判牍中的道德判断以及后出的叙述都须回到产生它们的场景，而不可直接代表全域动机或人口。

将事件与地方层次连接时，应同时考察资源怎样被征集、信息怎样传递、迁徙怎样改变户籍和劳动、宗教与亲属网络怎样提供或限制协作。现有证据往往只照亮少数地点；保留这种不均衡，才能避免把宋辽夏金元的多政权历史改写成单一中心的必然过程。

账本条目\texttt{NS-0976-REGNAL-YEAR}所记976年的“宋以宋太祖在位第17年作为诏令、赋役与外交文书的纪年框架”，属于宋的历史语境。条目所示前因是新岁颁历、年号和纪年是中枢与地方行政沟通的共同前提，其可见后果为为同年政令、战事和地方材料提供日期锚点，不能单独推知具体政策执行。这一变化若进入区域生活，必须通过道路、文书、军队、市场、家庭和地方中介才会发生；政权易手或法令发布不等于所有城镇、乡村与边地在同一时间获得同样经验。

本条所列来源为《宋史》卷1—23〈本纪〉；《续资治通鉴长编》本年条；《宋会要辑稿》相关门类。它们能够支持特定的年序、制度语言、政治行动或局部遗存，却不能无条件化为社会总量。账本特别提醒“纪年框架可靠；该记录仅确证政权纪年连续，年内具体事项须由本年条另证”。因此，军报的兵数、条约的名分、官府的族群分类、判牍中的道德判断以及后出的叙述都须回到产生它们的场景，而不可直接代表全域动机或人口。

将事件与地方层次连接时，应同时考察资源怎样被征集、信息怎样传递、迁徙怎样改变户籍和劳动、宗教与亲属网络怎样提供或限制协作。现有证据往往只照亮少数地点；保留这种不均衡，才能避免把宋辽夏金元的多政权历史改写成单一中心的必然过程。

账本条目\texttt{NS-0976-TAIZU-SUCCESSION}所记976年十月的“太祖去世，赵光义即位为太宗”，属于宋的历史语境。条目所示前因是太祖猝逝而皇子与弟弟的继承次序未在现存材料中完全透明，其可见后果为太宗获得皇位并调整北伐议程，继承合法性成为后世政治记忆的争点。这一变化若进入区域生活，必须通过道路、文书、军队、市场、家庭和地方中介才会发生；政权易手或法令发布不等于所有城镇、乡村与边地在同一时间获得同样经验。

本条所列来源为《宋史》卷2〈太祖本纪二〉；《宋史》卷3〈太宗本纪一〉；《续资治通鉴长编》开宝九年十月条。它们能够支持特定的年序、制度语言、政治行动或局部遗存，却不能无条件化为社会总量。账本特别提醒“即位日期可靠；烛影斧声为后起传闻，不能据以断言继承阴谋”。因此，军报的兵数、条约的名分、官府的族群分类、判牍中的道德判断以及后出的叙述都须回到产生它们的场景，而不可直接代表全域动机或人口。

将事件与地方层次连接时，应同时考察资源怎样被征集、信息怎样传递、迁徙怎样改变户籍和劳动、宗教与亲属网络怎样提供或限制协作。现有证据往往只照亮少数地点；保留这种不均衡，才能避免把宋辽夏金元的多政权历史改写成单一中心的必然过程。

账本条目\texttt{NS-0977-REGNAL-YEAR}所记977年的“宋以宋太宗在位第1年作为诏令、赋役与外交文书的纪年框架”，属于宋的历史语境。条目所示前因是新岁颁历、年号和纪年是中枢与地方行政沟通的共同前提，其可见后果为为同年政令、战事和地方材料提供日期锚点，不能单独推知具体政策执行。这一变化若进入区域生活，必须通过道路、文书、军队、市场、家庭和地方中介才会发生；政权易手或法令发布不等于所有城镇、乡村与边地在同一时间获得同样经验。

本条所列来源为《宋史》卷1—23〈本纪〉；《续资治通鉴长编》本年条；《宋会要辑稿》相关门类。它们能够支持特定的年序、制度语言、政治行动或局部遗存，却不能无条件化为社会总量。账本特别提醒“纪年框架可靠；该记录仅确证政权纪年连续，年内具体事项须由本年条另证”。因此，军报的兵数、条约的名分、官府的族群分类、判牍中的道德判断以及后出的叙述都须回到产生它们的场景，而不可直接代表全域动机或人口。

将事件与地方层次连接时，应同时考察资源怎样被征集、信息怎样传递、迁徙怎样改变户籍和劳动、宗教与亲属网络怎样提供或限制协作。现有证据往往只照亮少数地点；保留这种不均衡，才能避免把宋辽夏金元的多政权历史改写成单一中心的必然过程。

账本条目\texttt{NS-0978-REGNAL-YEAR}所记978年的“宋以宋太宗在位第2年作为诏令、赋役与外交文书的纪年框架”，属于宋的历史语境。条目所示前因是新岁颁历、年号和纪年是中枢与地方行政沟通的共同前提，其可见后果为为同年政令、战事和地方材料提供日期锚点，不能单独推知具体政策执行。这一变化若进入区域生活，必须通过道路、文书、军队、市场、家庭和地方中介才会发生；政权易手或法令发布不等于所有城镇、乡村与边地在同一时间获得同样经验。

本条所列来源为《宋史》卷1—23〈本纪〉；《续资治通鉴长编》本年条；《宋会要辑稿》相关门类。它们能够支持特定的年序、制度语言、政治行动或局部遗存，却不能无条件化为社会总量。账本特别提醒“纪年框架可靠；该记录仅确证政权纪年连续，年内具体事项须由本年条另证”。因此，军报的兵数、条约的名分、官府的族群分类、判牍中的道德判断以及后出的叙述都须回到产生它们的场景，而不可直接代表全域动机或人口。

将事件与地方层次连接时，应同时考察资源怎样被征集、信息怎样传递、迁徙怎样改变户籍和劳动、宗教与亲属网络怎样提供或限制协作。现有证据往往只照亮少数地点；保留这种不均衡，才能避免把宋辽夏金元的多政权历史改写成单一中心的必然过程。

账本条目\texttt{NS-0979-GAOLIANG-RIVER}所记979年六月的“宋军北攻幽州，在高梁河失利后南撤”，属于宋与辽的历史语境。条目所示前因是北汉既灭，太宗试图趁势夺取燕云；辽军可依近畿与骑兵机动应战，其可见后果为宋放弃速取燕云的设想，北部边防和对辽外交成为长期财政军事负担。这一变化若进入区域生活，必须通过道路、文书、军队、市场、家庭和地方中介才会发生；政权易手或法令发布不等于所有城镇、乡村与边地在同一时间获得同样经验。

本条所列来源为《宋史》卷3〈太宗本纪一〉；《辽史》〈景宗本纪〉；《续资治通鉴长编》太平兴国四年六月条。它们能够支持特定的年序、制度语言、政治行动或局部遗存，却不能无条件化为社会总量。账本特别提醒“战败可靠；宋辽双方兵力、皇帝受伤和追击路线说法不一”。因此，军报的兵数、条约的名分、官府的族群分类、判牍中的道德判断以及后出的叙述都须回到产生它们的场景，而不可直接代表全域动机或人口。

将事件与地方层次连接时，应同时考察资源怎样被征集、信息怎样传递、迁徙怎样改变户籍和劳动、宗教与亲属网络怎样提供或限制协作。现有证据往往只照亮少数地点；保留这种不均衡，才能避免把宋辽夏金元的多政权历史改写成单一中心的必然过程。

账本条目\texttt{NS-0979-NORTHERN-HAN}所记979年五月的“宋灭北汉，完成对中原和河东割据政权的统一”，属于宋与北汉的历史语境。条目所示前因是北汉长期依赖辽援，宋在统一南方后可集中兵力攻太原，其可见后果为河东并入宋朝，但宋与辽失去缓冲政权，边境冲突直接化。这一变化若进入区域生活，必须通过道路、文书、军队、市场、家庭和地方中介才会发生；政权易手或法令发布不等于所有城镇、乡村与边地在同一时间获得同样经验。

本条所列来源为《宋史》卷3〈太宗本纪一〉；《宋史》卷480〈北汉世家〉；《续资治通鉴长编》太平兴国四年五月条。它们能够支持特定的年序、制度语言、政治行动或局部遗存，却不能无条件化为社会总量。账本特别提醒“灭国事实可靠；辽援军作用与城内处境材料不均”。因此，军报的兵数、条约的名分、官府的族群分类、判牍中的道德判断以及后出的叙述都须回到产生它们的场景，而不可直接代表全域动机或人口。

将事件与地方层次连接时，应同时考察资源怎样被征集、信息怎样传递、迁徙怎样改变户籍和劳动、宗教与亲属网络怎样提供或限制协作。现有证据往往只照亮少数地点；保留这种不均衡，才能避免把宋辽夏金元的多政权历史改写成单一中心的必然过程。

账本条目\texttt{NS-0979-REGNAL-YEAR}所记979年的“宋以宋太宗在位第3年作为诏令、赋役与外交文书的纪年框架”，属于宋的历史语境。条目所示前因是新岁颁历、年号和纪年是中枢与地方行政沟通的共同前提，其可见后果为为同年政令、战事和地方材料提供日期锚点，不能单独推知具体政策执行。这一变化若进入区域生活，必须通过道路、文书、军队、市场、家庭和地方中介才会发生；政权易手或法令发布不等于所有城镇、乡村与边地在同一时间获得同样经验。

本条所列来源为《宋史》卷1—23〈本纪〉；《续资治通鉴长编》本年条；《宋会要辑稿》相关门类。它们能够支持特定的年序、制度语言、政治行动或局部遗存，却不能无条件化为社会总量。账本特别提醒“纪年框架可靠；该记录仅确证政权纪年连续，年内具体事项须由本年条另证”。因此，军报的兵数、条约的名分、官府的族群分类、判牍中的道德判断以及后出的叙述都须回到产生它们的场景，而不可直接代表全域动机或人口。

将事件与地方层次连接时，应同时考察资源怎样被征集、信息怎样传递、迁徙怎样改变户籍和劳动、宗教与亲属网络怎样提供或限制协作。现有证据往往只照亮少数地点；保留这种不均衡，才能避免把宋辽夏金元的多政权历史改写成单一中心的必然过程。

账本条目\texttt{NS-0980-REGNAL-YEAR}所记980年的“宋以宋太宗在位第4年作为诏令、赋役与外交文书的纪年框架”，属于宋的历史语境。条目所示前因是新岁颁历、年号和纪年是中枢与地方行政沟通的共同前提，其可见后果为为同年政令、战事和地方材料提供日期锚点，不能单独推知具体政策执行。这一变化若进入区域生活，必须通过道路、文书、军队、市场、家庭和地方中介才会发生；政权易手或法令发布不等于所有城镇、乡村与边地在同一时间获得同样经验。

本条所列来源为《宋史》卷1—23〈本纪〉；《续资治通鉴长编》本年条；《宋会要辑稿》相关门类。它们能够支持特定的年序、制度语言、政治行动或局部遗存，却不能无条件化为社会总量。账本特别提醒“纪年框架可靠；该记录仅确证政权纪年连续，年内具体事项须由本年条另证”。因此，军报的兵数、条约的名分、官府的族群分类、判牍中的道德判断以及后出的叙述都须回到产生它们的场景，而不可直接代表全域动机或人口。

将事件与地方层次连接时，应同时考察资源怎样被征集、信息怎样传递、迁徙怎样改变户籍和劳动、宗教与亲属网络怎样提供或限制协作。现有证据往往只照亮少数地点；保留这种不均衡，才能避免把宋辽夏金元的多政权历史改写成单一中心的必然过程。

账本条目\texttt{NS-0981-REGNAL-YEAR}所记981年的“宋以宋太宗在位第5年作为诏令、赋役与外交文书的纪年框架”，属于宋的历史语境。条目所示前因是新岁颁历、年号和纪年是中枢与地方行政沟通的共同前提，其可见后果为为同年政令、战事和地方材料提供日期锚点，不能单独推知具体政策执行。这一变化若进入区域生活，必须通过道路、文书、军队、市场、家庭和地方中介才会发生；政权易手或法令发布不等于所有城镇、乡村与边地在同一时间获得同样经验。

本条所列来源为《宋史》卷1—23〈本纪〉；《续资治通鉴长编》本年条；《宋会要辑稿》相关门类。它们能够支持特定的年序、制度语言、政治行动或局部遗存，却不能无条件化为社会总量。账本特别提醒“纪年框架可靠；该记录仅确证政权纪年连续，年内具体事项须由本年条另证”。因此，军报的兵数、条约的名分、官府的族群分类、判牍中的道德判断以及后出的叙述都须回到产生它们的场景，而不可直接代表全域动机或人口。

将事件与地方层次连接时，应同时考察资源怎样被征集、信息怎样传递、迁徙怎样改变户籍和劳动、宗教与亲属网络怎样提供或限制协作。现有证据往往只照亮少数地点；保留这种不均衡，才能避免把宋辽夏金元的多政权历史改写成单一中心的必然过程。

账本条目\texttt{NS-0982-REGNAL-YEAR}所记982年的“宋以宋太宗在位第6年作为诏令、赋役与外交文书的纪年框架”，属于宋的历史语境。条目所示前因是新岁颁历、年号和纪年是中枢与地方行政沟通的共同前提，其可见后果为为同年政令、战事和地方材料提供日期锚点，不能单独推知具体政策执行。这一变化若进入区域生活，必须通过道路、文书、军队、市场、家庭和地方中介才会发生；政权易手或法令发布不等于所有城镇、乡村与边地在同一时间获得同样经验。

本条所列来源为《宋史》卷1—23〈本纪〉；《续资治通鉴长编》本年条；《宋会要辑稿》相关门类。它们能够支持特定的年序、制度语言、政治行动或局部遗存，却不能无条件化为社会总量。账本特别提醒“纪年框架可靠；该记录仅确证政权纪年连续，年内具体事项须由本年条另证”。因此，军报的兵数、条约的名分、官府的族群分类、判牍中的道德判断以及后出的叙述都须回到产生它们的场景，而不可直接代表全域动机或人口。

将事件与地方层次连接时，应同时考察资源怎样被征集、信息怎样传递、迁徙怎样改变户籍和劳动、宗教与亲属网络怎样提供或限制协作。现有证据往往只照亮少数地点；保留这种不均衡，才能避免把宋辽夏金元的多政权历史改写成单一中心的必然过程。

账本条目\texttt{NS-0983-REGNAL-YEAR}所记983年的“宋以宋太宗在位第7年作为诏令、赋役与外交文书的纪年框架”，属于宋的历史语境。条目所示前因是新岁颁历、年号和纪年是中枢与地方行政沟通的共同前提，其可见后果为为同年政令、战事和地方材料提供日期锚点，不能单独推知具体政策执行。这一变化若进入区域生活，必须通过道路、文书、军队、市场、家庭和地方中介才会发生；政权易手或法令发布不等于所有城镇、乡村与边地在同一时间获得同样经验。

本条所列来源为《宋史》卷1—23〈本纪〉；《续资治通鉴长编》本年条；《宋会要辑稿》相关门类。它们能够支持特定的年序、制度语言、政治行动或局部遗存，却不能无条件化为社会总量。账本特别提醒“纪年框架可靠；该记录仅确证政权纪年连续，年内具体事项须由本年条另证”。因此，军报的兵数、条约的名分、官府的族群分类、判牍中的道德判断以及后出的叙述都须回到产生它们的场景，而不可直接代表全域动机或人口。

将事件与地方层次连接时，应同时考察资源怎样被征集、信息怎样传递、迁徙怎样改变户籍和劳动、宗教与亲属网络怎样提供或限制协作。现有证据往往只照亮少数地点；保留这种不均衡，才能避免把宋辽夏金元的多政权历史改写成单一中心的必然过程。

\subsection{宗教、法律与材料边界}

账本条目\texttt{NS-0984-REGNAL-YEAR}所记984年的“宋以宋太宗在位第8年作为诏令、赋役与外交文书的纪年框架”，属于宋的历史语境。条目所示前因是新岁颁历、年号和纪年是中枢与地方行政沟通的共同前提，其可见后果为为同年政令、战事和地方材料提供日期锚点，不能单独推知具体政策执行。这一变化若进入区域生活，必须通过道路、文书、军队、市场、家庭和地方中介才会发生；政权易手或法令发布不等于所有城镇、乡村与边地在同一时间获得同样经验。

本条所列来源为《宋史》卷1—23〈本纪〉；《续资治通鉴长编》本年条；《宋会要辑稿》相关门类。它们能够支持特定的年序、制度语言、政治行动或局部遗存，却不能无条件化为社会总量。账本特别提醒“纪年框架可靠；该记录仅确证政权纪年连续，年内具体事项须由本年条另证”。因此，军报的兵数、条约的名分、官府的族群分类、判牍中的道德判断以及后出的叙述都须回到产生它们的场景，而不可直接代表全域动机或人口。

将事件与地方层次连接时，应同时考察资源怎样被征集、信息怎样传递、迁徙怎样改变户籍和劳动、宗教与亲属网络怎样提供或限制协作。现有证据往往只照亮少数地点；保留这种不均衡，才能避免把宋辽夏金元的多政权历史改写成单一中心的必然过程。

账本条目\texttt{NS-0985-REGNAL-YEAR}所记985年的“宋以宋太宗在位第9年作为诏令、赋役与外交文书的纪年框架”，属于宋的历史语境。条目所示前因是新岁颁历、年号和纪年是中枢与地方行政沟通的共同前提，其可见后果为为同年政令、战事和地方材料提供日期锚点，不能单独推知具体政策执行。这一变化若进入区域生活，必须通过道路、文书、军队、市场、家庭和地方中介才会发生；政权易手或法令发布不等于所有城镇、乡村与边地在同一时间获得同样经验。

本条所列来源为《宋史》卷1—23〈本纪〉；《续资治通鉴长编》本年条；《宋会要辑稿》相关门类。它们能够支持特定的年序、制度语言、政治行动或局部遗存，却不能无条件化为社会总量。账本特别提醒“纪年框架可靠；该记录仅确证政权纪年连续，年内具体事项须由本年条另证”。因此，军报的兵数、条约的名分、官府的族群分类、判牍中的道德判断以及后出的叙述都须回到产生它们的场景，而不可直接代表全域动机或人口。

将事件与地方层次连接时，应同时考察资源怎样被征集、信息怎样传递、迁徙怎样改变户籍和劳动、宗教与亲属网络怎样提供或限制协作。现有证据往往只照亮少数地点；保留这种不均衡，才能避免把宋辽夏金元的多政权历史改写成单一中心的必然过程。

账本条目\texttt{NS-0986-REGNAL-YEAR}所记986年的“宋以宋太宗在位第10年作为诏令、赋役与外交文书的纪年框架”，属于宋的历史语境。条目所示前因是新岁颁历、年号和纪年是中枢与地方行政沟通的共同前提，其可见后果为为同年政令、战事和地方材料提供日期锚点，不能单独推知具体政策执行。这一变化若进入区域生活，必须通过道路、文书、军队、市场、家庭和地方中介才会发生；政权易手或法令发布不等于所有城镇、乡村与边地在同一时间获得同样经验。

本条所列来源为《宋史》卷1—23〈本纪〉；《续资治通鉴长编》本年条；《宋会要辑稿》相关门类。它们能够支持特定的年序、制度语言、政治行动或局部遗存，却不能无条件化为社会总量。账本特别提醒“纪年框架可靠；该记录仅确证政权纪年连续，年内具体事项须由本年条另证”。因此，军报的兵数、条约的名分、官府的族群分类、判牍中的道德判断以及后出的叙述都须回到产生它们的场景，而不可直接代表全域动机或人口。

将事件与地方层次连接时，应同时考察资源怎样被征集、信息怎样传递、迁徙怎样改变户籍和劳动、宗教与亲属网络怎样提供或限制协作。现有证据往往只照亮少数地点；保留这种不均衡，才能避免把宋辽夏金元的多政权历史改写成单一中心的必然过程。

账本条目\texttt{NS-0986-YONGXI-CAMPAIGN}所记986年的“宋以三路军北伐辽，随后遭挫撤军”，属于宋与辽的历史语境。条目所示前因是宋廷仍欲夺回燕云，设想以多路并进牵制辽军，其可见后果为失利强化守势边防与岁费压力，也影响宋廷对将领授权的方式。这一变化若进入区域生活，必须通过道路、文书、军队、市场、家庭和地方中介才会发生；政权易手或法令发布不等于所有城镇、乡村与边地在同一时间获得同样经验。

本条所列来源为《宋史》卷4〈太宗本纪二〉；《辽史》〈圣宗本纪〉；《续资治通鉴长编》雍熙三年条。它们能够支持特定的年序、制度语言、政治行动或局部遗存，却不能无条件化为社会总量。账本特别提醒“战役序列可靠；将领责任、军数及各路协同失败须对读双方史书”。因此，军报的兵数、条约的名分、官府的族群分类、判牍中的道德判断以及后出的叙述都须回到产生它们的场景，而不可直接代表全域动机或人口。

将事件与地方层次连接时，应同时考察资源怎样被征集、信息怎样传递、迁徙怎样改变户籍和劳动、宗教与亲属网络怎样提供或限制协作。现有证据往往只照亮少数地点；保留这种不均衡，才能避免把宋辽夏金元的多政权历史改写成单一中心的必然过程。

账本条目\texttt{NS-0987-REGNAL-YEAR}所记987年的“宋以宋太宗在位第11年作为诏令、赋役与外交文书的纪年框架”，属于宋的历史语境。条目所示前因是新岁颁历、年号和纪年是中枢与地方行政沟通的共同前提，其可见后果为为同年政令、战事和地方材料提供日期锚点，不能单独推知具体政策执行。这一变化若进入区域生活，必须通过道路、文书、军队、市场、家庭和地方中介才会发生；政权易手或法令发布不等于所有城镇、乡村与边地在同一时间获得同样经验。

本条所列来源为《宋史》卷1—23〈本纪〉；《续资治通鉴长编》本年条；《宋会要辑稿》相关门类。它们能够支持特定的年序、制度语言、政治行动或局部遗存，却不能无条件化为社会总量。账本特别提醒“纪年框架可靠；该记录仅确证政权纪年连续，年内具体事项须由本年条另证”。因此，军报的兵数、条约的名分、官府的族群分类、判牍中的道德判断以及后出的叙述都须回到产生它们的场景，而不可直接代表全域动机或人口。

将事件与地方层次连接时，应同时考察资源怎样被征集、信息怎样传递、迁徙怎样改变户籍和劳动、宗教与亲属网络怎样提供或限制协作。现有证据往往只照亮少数地点；保留这种不均衡，才能避免把宋辽夏金元的多政权历史改写成单一中心的必然过程。

账本条目\texttt{NS-0988-REGNAL-YEAR}所记988年的“宋以宋太宗在位第12年作为诏令、赋役与外交文书的纪年框架”，属于宋的历史语境。条目所示前因是新岁颁历、年号和纪年是中枢与地方行政沟通的共同前提，其可见后果为为同年政令、战事和地方材料提供日期锚点，不能单独推知具体政策执行。这一变化若进入区域生活，必须通过道路、文书、军队、市场、家庭和地方中介才会发生；政权易手或法令发布不等于所有城镇、乡村与边地在同一时间获得同样经验。

本条所列来源为《宋史》卷1—23〈本纪〉；《续资治通鉴长编》本年条；《宋会要辑稿》相关门类。它们能够支持特定的年序、制度语言、政治行动或局部遗存，却不能无条件化为社会总量。账本特别提醒“纪年框架可靠；该记录仅确证政权纪年连续，年内具体事项须由本年条另证”。因此，军报的兵数、条约的名分、官府的族群分类、判牍中的道德判断以及后出的叙述都须回到产生它们的场景，而不可直接代表全域动机或人口。

将事件与地方层次连接时，应同时考察资源怎样被征集、信息怎样传递、迁徙怎样改变户籍和劳动、宗教与亲属网络怎样提供或限制协作。现有证据往往只照亮少数地点；保留这种不均衡，才能避免把宋辽夏金元的多政权历史改写成单一中心的必然过程。

账本条目\texttt{NS-0989-REGNAL-YEAR}所记989年的“宋以宋太宗在位第13年作为诏令、赋役与外交文书的纪年框架”，属于宋的历史语境。条目所示前因是新岁颁历、年号和纪年是中枢与地方行政沟通的共同前提，其可见后果为为同年政令、战事和地方材料提供日期锚点，不能单独推知具体政策执行。这一变化若进入区域生活，必须通过道路、文书、军队、市场、家庭和地方中介才会发生；政权易手或法令发布不等于所有城镇、乡村与边地在同一时间获得同样经验。

本条所列来源为《宋史》卷1—23〈本纪〉；《续资治通鉴长编》本年条；《宋会要辑稿》相关门类。它们能够支持特定的年序、制度语言、政治行动或局部遗存，却不能无条件化为社会总量。账本特别提醒“纪年框架可靠；该记录仅确证政权纪年连续，年内具体事项须由本年条另证”。因此，军报的兵数、条约的名分、官府的族群分类、判牍中的道德判断以及后出的叙述都须回到产生它们的场景，而不可直接代表全域动机或人口。

将事件与地方层次连接时，应同时考察资源怎样被征集、信息怎样传递、迁徙怎样改变户籍和劳动、宗教与亲属网络怎样提供或限制协作。现有证据往往只照亮少数地点；保留这种不均衡，才能避免把宋辽夏金元的多政权历史改写成单一中心的必然过程。

账本条目\texttt{NS-0990-REGNAL-YEAR}所记990年的“宋以宋太宗在位第14年作为诏令、赋役与外交文书的纪年框架”，属于宋的历史语境。条目所示前因是新岁颁历、年号和纪年是中枢与地方行政沟通的共同前提，其可见后果为为同年政令、战事和地方材料提供日期锚点，不能单独推知具体政策执行。这一变化若进入区域生活，必须通过道路、文书、军队、市场、家庭和地方中介才会发生；政权易手或法令发布不等于所有城镇、乡村与边地在同一时间获得同样经验。

本条所列来源为《宋史》卷1—23〈本纪〉；《续资治通鉴长编》本年条；《宋会要辑稿》相关门类。它们能够支持特定的年序、制度语言、政治行动或局部遗存，却不能无条件化为社会总量。账本特别提醒“纪年框架可靠；该记录仅确证政权纪年连续，年内具体事项须由本年条另证”。因此，军报的兵数、条约的名分、官府的族群分类、判牍中的道德判断以及后出的叙述都须回到产生它们的场景，而不可直接代表全域动机或人口。

将事件与地方层次连接时，应同时考察资源怎样被征集、信息怎样传递、迁徙怎样改变户籍和劳动、宗教与亲属网络怎样提供或限制协作。现有证据往往只照亮少数地点；保留这种不均衡，才能避免把宋辽夏金元的多政权历史改写成单一中心的必然过程。

账本条目\texttt{NS-0991-REGNAL-YEAR}所记991年的“宋以宋太宗在位第15年作为诏令、赋役与外交文书的纪年框架”，属于宋的历史语境。条目所示前因是新岁颁历、年号和纪年是中枢与地方行政沟通的共同前提，其可见后果为为同年政令、战事和地方材料提供日期锚点，不能单独推知具体政策执行。这一变化若进入区域生活，必须通过道路、文书、军队、市场、家庭和地方中介才会发生；政权易手或法令发布不等于所有城镇、乡村与边地在同一时间获得同样经验。

本条所列来源为《宋史》卷1—23〈本纪〉；《续资治通鉴长编》本年条；《宋会要辑稿》相关门类。它们能够支持特定的年序、制度语言、政治行动或局部遗存，却不能无条件化为社会总量。账本特别提醒“纪年框架可靠；该记录仅确证政权纪年连续，年内具体事项须由本年条另证”。因此，军报的兵数、条约的名分、官府的族群分类、判牍中的道德判断以及后出的叙述都须回到产生它们的场景，而不可直接代表全域动机或人口。

将事件与地方层次连接时，应同时考察资源怎样被征集、信息怎样传递、迁徙怎样改变户籍和劳动、宗教与亲属网络怎样提供或限制协作。现有证据往往只照亮少数地点；保留这种不均衡，才能避免把宋辽夏金元的多政权历史改写成单一中心的必然过程。
